\pdfbookmark[0]{Abstract}{Abstract}
\chapter*{Abstract}
\markboth{Abstract}{Abstract}

This thesis addresses the advanced 2D and 3D characterisation of clinker phases and hydrated cementitious binders.
Conventional imaging techniques often do not provide the required combination of spatial resolution and suitable signals, such as high chemical sensitivity or information on mechanical properties, to correlate these spatially resolved information with each other or with higher-resolution methods such as provided by \acl*{SEM} (\acs*{SEM}).
An objective of this work is to overcome these limitations by establishing correlative microscopy workflows and automated evaluation tools that bridge the gap between high-resolution imaging and physical property determination.
To achieve this, various modern imaging and analytical techniques, including \acs*{SEM}, focussed ion beam nano tomography (\acs*{FIBnt}), \acl*{LAICPMS} (\acs*{LAICPMS}), \acl*{ArBIB} (\acs*{ArBIB}), \acl*{LTDSC} (\acs*{LTDSC}), and \acl*{NI} (\acs*{NI}), are systematically combined.

A central component of this research is the use of advanced preparation techniques.
\acs*{ArBIB} polishing is employed to create surfaces with very few artefacts, which is crucial for high-resolution \acs*{SEM} imaging as it preserves the native pore structure without the need for resin embedding.
This method proved essential for preparing samples for subsequent correlative analyses.

The advantages of the correlative approach are demonstrated through several workflows.
The combination of \acs*{SEM}-\acl*{EDX} with \acs*{NI} or \acs*{LAICPMS} allows for the direct assignment of phase-specific mechanical properties and trace element distributions.
This moves beyond traditional deconvolution methods by linking chemical and mechanical data to specific microstructural features, such as assigning trace elements to specific clinker phase or determining the hardness of individual hydrate phases.

To move from qualitative observation of microstructural features to quantitative analysis, automated image analysis tools were developed.
A novel algorithm enables the automated measurement of the \acl*{HLT} (\acs*{HLT}) and inter-particle distances from large-scale 2D images.
This facilitates the statistical analysis of hydration kinetics, as demonstrated on \acs*{C3S} and \acs*{C2S} systems, and quantifies the dispersion effects of admixtures or physical treatments like \acl*{PUS}.
Furthermore, the correlation of the \acs*{HLT} measurements with the crystal orientation measured by \acl*{EBSD} of \ce{C3S} indicates that \acs*{CSH} growth around \acs*{C3S} particles is governed more by the local availability of space for crystallisation in the pores than by the crystallographic orientation of the original clinker grain (\acs*{C3S}).

For the 3D characterisation of hydrated cements, optimised workflows for \acs*{FIBnt} were established.
An improved lift-out technique for cementitious materials was implemented to minimise imaging artefacts, and \acl*{fs-laser} preparation was utilised to rapidly access larger, more representative volumes.
These preparation methods are complemented by an automated post-processing pipeline for 3D data, which includes registration, denoising, and superpixel-based segmentation to generate quantitative phase maps.
These 3D datasets, alongside 2D imaging, \acs*{LTDSC}, and \acl*{MIP}, were used in a comparative study of gel and capillary porosity, revealing a linear correlation between the reduction of capillary porosity and the formation of gel porosity during hydration.

The results demonstrate that a correlative, multi-method approach yields a more comprehensive understanding than individual techniques can provide alone.
The developed automated tools provide objective and statistically robust data, laying the foundation for future research aimed at optimising binder efficiency and durability through a deeper, data-driven understanding of cementitious microstructures.